% TOP_PROBLEM: {"problemId": "problem.pisot-substitution-conjecture", "canonicalTitle": "Pisot Substitution Conjecture", "releaseRank": 462, "primaryDomain": "probability_ergodic_dynamics", "primaryDomainLabel": "Probability, ergodic theory & dynamics", "displayStatus": "open_with_solved_subcases", "releaseStatus": "open_with_solved_subcases"}
% !TeX program = lualatex
% ATTEMPT_STATUS: unresolved
\documentclass[11pt]{article}
\usepackage[margin=25mm]{geometry}
\usepackage{fontspec}
\setmainfont{Latin Modern Roman}
\usepackage{amsmath,amssymb,mathtools}
\usepackage{unicode-math}
\setmathfont{Latin Modern Math}
\usepackage{xurl,hyperref}
\hypersetup{colorlinks=true,urlcolor=blue}
\setcounter{secnumdepth}{0}
\setlength{\parindent}{0pt}
\setlength{\parskip}{5pt}
\setlength{\emergencystretch}{3em}
\begin{document}
\section*{\texorpdfstring{Pisot Substitution Conjecture}{Pisot Substitution Conjecture}}\label{462}
\subsection{Definitions and mathematical statement}
A substitution assigns to each letter in a finite alphabet a nonempty word and extends by concatenation. Its incidence matrix $M$ has entry $M_{ij}$ equal to the number of occurrences of letter $i$ in the image of letter $j$. In the irreducible unimodular Pisot setting, the characteristic polynomial is irreducible over ${\mathbb{Q}}$, $\det M=\pm1$, and its dominant eigenvalue $\beta>1$ is a Pisot number: every other algebraic conjugate has modulus less than one.

The substitution subshift consists of bi-infinite words whose finite blocks occur in iterated substituted words, with the left shift and its invariant probability measure in the primitive setting of the source. Pure discrete spectrum means that the eigenfunctions of the Koopman operator $Uf=f\circ\sigma$ span $L^2$. The conjecture asserts this property for every substitution in the specified class. A symbolic periodicity statement or a pure-point diffraction assertion without the needed dynamical equivalence is not substituted for it.
\par\smallskip\textit{Formulation and concept references:} \href{https://ems.press/content/serial-article-files/32990?nt=1}{[S1]}.

\subsection{Short English statement}
Must every irreducible unimodular Pisot substitution system have a complete basis of measurable eigenfunctions?
\subsection{Research attempt}
\paragraph{Formulation and source audit (27 September 2026).}
The target is the measurable spectrum of the symbolic shift, with its
invariant probability measure. It is not a claim that the sequence is
periodic. An irrational rotation has infinitely many distinct orbit points
and nevertheless has pure discrete spectrum. Conversely producing some
eigenfunctions does not prove that they span all of $L^2$.

There is a bibliographic mismatch in the inherited source entry. Its EMS
URL resolves to Berth\'e--Steiner--Thuswaldner, \emph{Multidimensional
continued fractions and symbolic codings of toral translations}, JEMS 25
(2023), 4997--5057, rather than the title printed in [S1]. The PDF does
discuss the selected Pisot conjecture, notably Section 2.2. The entry is
preserved and the actual document identified in [E1]. Hollander--Solomyak
[E2] prove the two-symbol case. These statements do not extend that
theorem to an arbitrary alphabet. The special example below is proved
directly through a rotation coding, rather than by quoting its membership
in that solved class.

The attack has two levels. Pisot contraction gives a uniform discrepancy
bound for letter counts. In a worked Fibonacci example, one can additionally
construct an almost-everywhere one-to-one rotation coding, which gives
the full spectral conclusion. The general missing step is the analogue
of that one-to-one assertion; contraction alone does not prove it.

\paragraph{The matrix hypotheses really imply primitivity here.}
The incidence matrix is a nonnegative integer matrix. If its directed
graph were not strongly connected, a simultaneous permutation of rows
and columns would put it in proper block triangular form. Its characteristic
polynomial would then factor over $\mathbb Q$, contrary to irreducibility.
Thus $M$ is irreducible as a nonnegative matrix.

If its period were $h>1$, the cyclic decomposition of that graph would
give the eigenvalues $\beta\zeta$ for all $h$th roots of unity $\zeta$.
For example, multiply coordinates in successive cyclic classes by successive
powers of $\zeta$ to transform a Perron eigenvector into an eigenvector
with eigenvalue $\beta\zeta$ or its inverse-phase version. At least one
other eigenvalue would have modulus $\beta>1$. But irreducibility of the
characteristic polynomial makes all its roots conjugates of $\beta$,
and the Pisot hypothesis puts every other root strictly inside the unit
circle. Therefore $h=1$ and $M$ is primitive. The standard Perron--Frobenius
theorem consequently supplies positive left and right eigenvectors.

This argument uses the irreducible characteristic polynomial, not merely
the existence of a Pisot Perron eigenvalue in an otherwise arbitrary
substitution. Additional invariant blocks or additional eigenvalues are
not allowed in the selected formulation.

\paragraph{Lemma 462.1: bounded letter discrepancy from contraction.}
For a word $w$, let $\ell(w)\in\mathbb Z^d$ be its vector of letter
counts, and write $|w|=\mathbf1^{\mathsf T}\ell(w)$. There is a positive
vector $f$ with $\mathbf1^{\mathsf T}f=1$ and $Mf=\beta f$, and a
constant $C$ depending on the substitution, such that every word in its
language satisfies
\begin{equation}\label{a462:balance}
 \|\ell(w)-|w|f\|\le C.
\end{equation}

\emph{Proof.} Choose the positive left eigenvector $u^{\mathsf T}$ so
that $u^{\mathsf T}f=1$, and define the spectral projection
$Q=I-fu^{\mathsf T}$. It commutes with $M$ and annihilates the Perron
direction. Because the irreducible polynomial is separable in characteristic
zero, $M$ is diagonalizable over $\mathbb C$; all eigenvalues on the image
of $Q$ have modulus less than one. Thus, for some $0<r<1$ and $A<\infty$,
\[
 \|M^jQ\|\le Ar^j\quad(j\ge0).
\]

Let $p$ be a proper prefix of $\sigma^n(a)$. Inspect its cut at the
outermost substitution level. It consists of complete images under
$\sigma^{n-1}$ of a proper prefix of $\sigma(a)$, followed by a prefix
inside one remaining substituted letter. Iterating this decomposition gives
\[
 \ell(p)=\sum_{j=0}^{n-1}M^j\ell(p_j),
\]
where each $p_j$ is a proper prefix of the image of a single letter;
empty prefixes are permitted. There are finitely many possibilities for
these prefix vectors. If $B$ bounds their norms, then
\[
 \|Q\ell(p)\|\le AB\sum_{j\ge0}r^j=AB/(1-r).
\]
For a full prefix $p=\sigma^n(a)$, instead use
$Q\ell(p)=M^nQe_a$, which is uniformly bounded too. Every legal factor
is the difference of two prefixes of some $\sigma^n(a)$, so
$\|Q\ell(w)\|$ is uniformly bounded for all legal $w$.

It remains to use ordinary word length correctly. The projection $Q$
uses the left Perron vector $u$, which need not be $\mathbf1$.
Writing $z=Q\ell(w)$ gives
$\ell(w)=f(u^{\mathsf T}\ell(w))+z$ and
$|w|=u^{\mathsf T}\ell(w)+\mathbf1^{\mathsf T}z$.
Consequently
\[
 \ell(w)-|w|f=z-f\mathbf1^{\mathsf T}z.
\]
The preceding uniform bound on $z$ proves \eqref{a462:balance}.
\hfill$\square$

In particular two legal words of the same length have letter-count
vectors differing by at most $2C$. Also $f$ is their common limiting
letter-frequency vector. These are substantial combinatorial restrictions,
but the proof has only controlled $d$ counting observables. It has not
constructed a complete orthonormal family of eigenfunctions for all
measurable observables of the subshift.

\paragraph{A fully worked example: Fibonacci substitution.}
Take
\[
 \sigma(a)=ab,\qquad\sigma(b)=a,
 \qquad M=\begin{pmatrix}1&1\\1&0\end{pmatrix}.
\]
Its characteristic polynomial $z^2-z-1$ is irreducible, its determinant
is $-1$, and its roots are $\varphi=(1+\sqrt5)/2$ and $-1/\varphi$.
Thus it lies in the selected class. Set $\alpha=1/\varphi^2$.
We will identify its symbolic system with rotation by $\alpha$ on the
circle, up to sets of measure zero.

First define a one-sided word whose $b$'s occur at the zero-based positions
\begin{equation}\label{a462:beatty}
 B_k=\lfloor k\varphi^2\rfloor-1,\qquad k\ge1,
\end{equation}
and whose other letters are $a$. Its $a$-positions are
$A_k=\lfloor k\varphi\rfloor-1$. Here is a direct proof of that
partition. For every positive integer $m$, the numbers in the two Beatty
sequences $\lfloor k\varphi\rfloor$ and $\lfloor k\varphi^2\rfloor$
that are at most $m$ have total count
\[
 \left\lfloor\frac{m+1}{\varphi}\right\rfloor+
 \left\lfloor\frac{m+1}{\varphi^2}\right\rfloor=m,
\]
since $1/\varphi+1/\varphi^2=1$ and neither summand before taking
floors is integral. Subtracting the same identity at $m-1$ shows that
exactly one sequence contains each positive integer $m$.

Under substitution, the $k$th $a$ creates the $k$th $b$. Before that $a$
there are $k-1$ letters $a$ and $\lfloor k\varphi\rfloor-k$ letters
$b$. Their images have total length
$2(k-1)+\lfloor k\varphi\rfloor-k$. The new $b$, second in the image
of the $k$th $a$, is therefore at position
\[
 \lfloor k\varphi\rfloor+k-1
       =\lfloor k\varphi^2\rfloor-1=B_k.
\]
The substituted word has exactly the same $b$-positions, so it is fixed
by $\sigma$. It starts with $a$ and hence agrees with the increasing
prefix limit of $\sigma^n(a)$: every fixed word beginning with $a$
must begin with each of those prefixes, whose lengths tend to infinity.

\paragraph{From the exact word to a measure-theoretic rotation.}
The number of $b$'s in the first $n$ letters of \eqref{a462:beatty} is
$\lfloor(n+1)\alpha\rfloor$. Thus its letter at position $n$ is $b$
exactly when
\[
 \lfloor(n+2)\alpha\rfloor-\lfloor(n+1)\alpha\rfloor=1.
\]
For rotation $T(x)=x+\alpha\pmod1$, partition the circle into
$I_a=[0,1-\alpha)$ and $I_b=[1-\alpha,1)$. The displayed word is
the forward itinerary of $x=\alpha$ under this partition.

For a general $x$, code all iterates $T^nx$, $n\in\mathbb Z$. The set
of $x$ whose orbit meets a partition endpoint is countable and has
Lebesgue measure zero. Away from that set the coding commutes with
rotation and the shift. Every finite cylinder condition pulls back to
an intersection of finitely many translated intervals, with finitely
many boundary points. If it has nonempty interior, the dense orbit of
$\alpha$ visits that interior, so its word occurs in the fixed word.
Conversely a block of the fixed word is realized on a one-sided interval
next to its coding point, even if that point is an endpoint. Hence it
has positive Lebesgue measure and belongs to the rotation language.
The subshift obtained by closing the shift orbit has exactly this language.

For completeness, the fixed-word language equals the substitution language.
Every fixed-word block lies in a long enough $\sigma^n(a)$. Conversely,
$b$ occurs in $\sigma(a)$, so $\sigma^n(b)$ occurs in
$\sigma^{n+1}(a)$; all substituted-letter blocks therefore occur in the
fixed word. No unrelated coding system has been inserted.

The coding is one-to-one away from the endpoint orbits. To see this,
take distinct points $x,y$ and put $d=y-x\pmod1\ne0$. The proper
interval $I_b$ and its translate $I_b-d$ have symmetric difference
with nonempty interior: a proper single interval cannot be invariant
under a nonzero circle translation. The dense orbit of $x$ under
$T$ meets that interior, so for some $n$ exactly one of $T^nx,T^ny$
belongs to $I_b$. Their itineraries differ. Thus the coding is injective
off the countable orbit of the endpoints. It is a Borel map between
standard Borel spaces, so its inverse on this full-measure image is
measurable.

Irrational rotation has unique invariant probability measure, Lebesgue
measure: if a measure is invariant, its $k$th Fourier coefficient equals
$e^{2\pi i k\alpha}$ times itself and hence is zero for $k\ne0$;
density of trigonometric polynomials determines the measure. Uniform
frequencies for interval cylinders follow from the geometric-sum proof of
equidistribution and approximation of interval indicators by continuous
functions. They identify the invariant measure of the coded subshift
with the pushforward of Lebesgue measure. Thus the almost-everywhere
bijection above is a measure-theoretic conjugacy for the measure in the
problem statement.

\paragraph{The spectral conclusion in this subcase.}
On the circle, the functions $e_k(x)=e^{2\pi i kx}$, $k\in\mathbb Z$,
form an orthonormal basis of $L^2$ and satisfy
\[
 e_k(Tx)=e^{2\pi i k\alpha}e_k(x).
\]
Transport them through the measurable inverse coding. They remain a
complete orthonormal basis of eigenfunctions on the Fibonacci subshift.
This proves pure discrete spectrum in the precise symbolic sense of the
catalogue. It does not merely establish pure-point diffraction or a
property of a suspension flow. It also exhibits why the sequence need
not be periodic: $\alpha$ is irrational.

The proof reconstructs a classical example, with no novelty claim. The
general two-symbol theorem [E2] covers many additional substitutions;
the elementary Beatty argument here is not asserted to cover them all.

\paragraph{Where the general geometric route stops.}
Lemma 462.1 puts projected prefix vectors in a bounded part of the
contracting space. Taking their closure suggests compact geometric pieces
and a torus factor. Establishing a factor is not enough: the factor map
could identify several symbolic configurations on a set of positive
measure. Eigenfunctions pulled back from the factor then span only the
functions constant on its fibers, rather than all of $L^2$.

The logical issue is visible without substitution terminology. The product
of an irrational circle rotation with a two-sided independent binary shift
has the circle rotation as a factor. That factor has a complete Fourier
basis, but the product system is not pure discrete. Indeed the centered
coordinate function $F(x,\omega)=\omega_0-1/2$ has
$\langle U^nF,F\rangle=0$ for every nonzero integer $n$. If a nonzero
function belonged to the pure-point subspace, expanding it in an
orthonormal eigenbasis would make its correlations
$\sum_j c_j\lambda_j^n$ with $c_j\ge0$ and $\sum c_j>0$; their
Cesaro mean squared modulus has a positive contribution from at least
one eigenvalue and cannot vanish. This contradicts the displayed
correlations. The product is not claimed to be a Pisot substitution;
it is a precise counterexample to the abstract ``pure-point factor implies
pure-point extension'' shortcut.

In the Fibonacci proof, separation by interval translates supplied almost-everywhere
injectivity. In higher-dimensional contracting spaces that role requires
control of overlaps and multiplicities of geometric pieces. Neither a
bounded projected orbit nor a numerical picture with small apparent
overlap certifies that overlap has measure zero. An algorithm checking
one substitution also does not establish successful termination with the
necessary coincidence property for every substitution in the class.

\paragraph{Finite checks and outcome.}
The companion computation compares a long exact Fibonacci fixed-word
prefix with the two Beatty descriptions, checks the substitution matrix
and its polynomial identity, and verifies the letter-count discrepancy
on all factors of bounded length in the generated word. The Beatty floors
are computed with integer square-root comparisons, not rounded floating
point values near integer boundaries. Finite agreement is a diagnostic;
the preceding identities prove the infinite coding.

\textbf{Outcome: UNRESOLVED, with a complete classical subcase.} The
attempt proves bounded letter discrepancy throughout the selected class
and proves the full measurable spectral assertion for the Fibonacci
substitution. The unresolved step for the general class is to turn
contraction and the resulting geometric factor into a sufficiently
one-to-one measurable model, or otherwise establish completeness of the
eigenfunctions. No general coincidence theorem is proved here.

\subsection{Sources}
\begin{itemize}
\item[S1]Substitutions, Rauzy fractals and tilings. Submitted scholarly source, accessed 2026-08-31.
\url{https://ems.press/content/serial-article-files/32990?nt=1}.
\textit{Formulation source.}
\item[E1]Val\'erie Berth\'e, Wolfgang Steiner, and J\"org M. Thuswaldner,
\emph{Multidimensional continued fractions and symbolic codings of toral
translations}, Journal of the European Mathematical Society 25 (2023),
4997--5057, Section 2.2.
\url{https://ems.press/content/serial-article-files/32990?nt=1}.
\textit{Actual primary document at the inherited URL, inspected
27 September 2026; the inherited bibliographic mismatch is recorded above.}
\item[E2]Michael Hollander and Boris Solomyak,
\emph{Two-symbol Pisot substitutions have pure discrete spectrum},
Ergodic Theory and Dynamical Systems 23 (2003), 533--540.
\url{https://doi.org/10.1017/S0143385702001384}.
\textit{Primary publisher abstract inspected 27 September 2026;
the two-symbol theorem is attributed, not generalized here.}
\end{itemize}
\end{document}
